% П2.4-1 — Равнобедренный треугольник: медиана к основанию
\documentclass[tikz,border=3pt]{standalone}
% Шапка всех рисунков: те же шрифты и стиль, что в книге.
\def\pfFontPath{../../шрифты/}
\usepackage{fontspec}
\usepackage{unicode-math}
\setmainfont{Manrope}[Path=\pfFontPath, Extension=.ttf,
  UprightFont=*-Regular, BoldFont=*-Bold,
  ItalicFont=*-Regular, ItalicFeatures={FakeSlant=0.18},
  BoldItalicFont=*-Bold, BoldItalicFeatures={FakeSlant=0.18}]
\setmathfont{FiraMath-Regular}[Path=\pfFontPath, Extension=.otf]
\usepackage{стиль/рисунки}

\begin{document}
\begin{tikzpicture}[pf, scale=0.9]
  \coordinate (A) at (0,0);
  \coordinate (C) at (5,0);
  \coordinate (B) at (2.5,3.6);
  \coordinate (M) at ($(A)!0.5!(C)$);   % середина AC — считается сама

  \draw[given, tick1] (A) -- (B);
  \draw[given, tick1] (B) -- (C);
  \draw[given, tick2] (A) -- (M);
  \draw[given, tick2] (M) -- (C);
  \draw[main] (B) -- (M);

  % равные углы при основании и равные половины угла B
  \pic[ang1] {angle = C--A--B};
  \pic[ang1] {angle = B--C--A};
  \pfAngTwo{A}{B}{M}{6.5}
  \pfAngTwo{M}{B}{C}{6.5}
  \pic[rightang] {right angle = B--M--C};

  \pfPoint{A}{A}{below left}
  \pfPoint{C}{C}{below right}
  \pfPoint{B}{B}{above}
  \pfPoint{M}{M}{below=4pt}
\end{tikzpicture}
\end{document}
